% Paper 2 · §6 Limitations and threats to validity
% Status: draft · 2026-05-05

\section{Limitations and threats to validity}
\label{sec:limitations}

\paragraph{Self-judging bias.}
We use DeepSeek V3.2 as both the answer-generating subject and the
LLM-as-judge. Self-judging is a known concern \citep{zheng2024judging}.
We mitigate by adopting the LongMemEval-S paper's fixed grading rubric
verbatim, and by cross-validating a 50-question stratified sample
with Gemini-2.5-pro acting as judge: absolute disagreement was
\(<\)1.5 pts. Full cross-judge replication on n=500 with a third model
is on our roadmap (estimated \cny{}30 additional cost).

\paragraph{Closed-haystack assumption.}
LongMemEval-S provides 50K-token closed haystacks where every relevant
session is in the corpus. Our results may not transfer to open-domain RAG
where relevant content must be identified among millions of documents.
In particular, our negative finding on Neo4j graph reranking
(\S\ref{sec:negative}.a) is closed-haystack-specific; we expect the
opposite tradeoff in open-domain settings (the original Zep result). We
plan a mixed-domain replication with NQ + LongMemEval-S as a future
benchmark.

\paragraph{Region-specific cost economics.}
Our cost analysis (\$3.50 USD per 500-question run) is specific to
Chinese cloud providers (Tencent Cloud T4 spot, Volc Ark coding plan).
Equivalent setup with AWS T4 + Anthropic Claude or OpenAI GPT-4o would
cost \(\geq\)\$50 per run. The 1/15 cost claim should be read as a
\emph{regional advantage}, not a fundamental architecture advantage.
That said, our pipeline is provider-neutral: replacing the Volc Ark
endpoint with any Anthropic-compatible API (Claude, Gemini, etc.) is
a one-line change.

\paragraph{Model availability.}
The Volc Ark coding plan provides nine open-via-API models at flat
rates that are unusual in the market. If this rate structure changes
or the plan is withdrawn, our cost claim degrades by \(5\)--\(15\times\).
We mitigate by maintaining provider-neutral SDK code and by mirroring
the bge-m3 + reranker weights to a self-hostable replication
(HuggingFace mirror at \texttt{hf-mirror.com}).

\paragraph{LongMemEval-S format constraints.}
LongMemEval-S evaluates 6 specific question types over chat
transcripts. Real-world memory tasks include task-completion
verification, multi-modal recall, and code understanding, which our
benchmark does not test. Compass v0.8's drift detection
(\S\ref{sec:method}.A) addresses the verification gap, but we do not
quantitatively benchmark it in this paper.

\paragraph{No live deployment data.}
All 500-question runs are offline reproductions. The companion plugin
runs in production at \url{https://compass.nautilus.social} since
2026-05-01 with a small user base; we do not yet report user-derived
accuracy because the user count and question diversity are too small
for statistical claims. We expect to report user-derived metrics in
a follow-up after \(\geq\)1000 active users.
